%-*-latex-*-

\documentclass[a4paper,11pt,twocolumn]{article}

% Language and fonts
%
\usepackage[utf8]{inputenc}
\usepackage[T1]{fontenc}
\usepackage[french]{babel}

\input{dnmacs.tex}


\begin{document}
\begin{center} {\Large Micro-dictionnaire bouddhique} \end{center}
\bigskip

{\dn aEtyogyAn} {\em atiyogay\={a}na} [pf. {\em ati} + {\em yoga} +
{\em y\={a}na}] n. vaj. «~Véhicule du yoga extraordinaire\fg,
dernier des six {\em y\={a}na} du Vajray\={a}na (ou dernier des trois
{\em y\={a}na} supérieurs). Discipline majeure de l'ordre tibétain
Nyingma (qui la dénomme aussi {\em Dzogchen}), elle est considérée
comme l'aboutissement des enseignements tantriques du bouddha
\'{S}\={a}kyamuni. \\

{\dn anA(m\qq{n}} {\em an\={a}tman} [pf. {\em an} + {\em \={a}tman}] m. bd.
«~Non-ego, non-noumène, absence d'existence en soi\fg. Ce n'est pas le
non-être (de Parménide) car la logique bouddhique ne reconnaît pas de
valeur ontologique au principe du tiers exclu. L'{\em an\={a}tman}
caractérise le bouddhisme et en constitue le pilier central. Sa logique
démontre que l'ego n'existe pas intrinsèquement car il ne
peut être trouvé ni dans les agrégats ({\em skandha}), c'est-à-dire
dans les parties issues de son analyse, ni dans leur union, résultante
de la synthèse. Ainsi le Ny\={a}natiloka dit: «~Notre existence
dite individuelle n'est en réalité rien de plus qu'un processus de
phénomènes physiques et psychiques, processus qui était en marche bien
avant notre naissance, depuis des temps immémoriaux, et qui se
poursuivra au-delà de notre mort pour une durée infiniment
longue. Qu'ils soient pris séparément ou tous ensemble, jamais ces
cinq agrégats ne constituent une unité individuelle et autonome
réelle; en dehors d'eux, il n'existe rien non plus que l'on puisse
désigner comme un Moi indépendant; la croyance en une entité
personnelle réelle, un Moi au sens suprême du terme, est pure
illusion.~» Le bouddhisme ainsi ne nie que l'existence
permanente en indépendance des agrégats (par exemple l'âme chrétienne
ou  l'{\em \={a}tman} hindouiste), et ne constitue pas un
nihilisme (au sens d'adhésion au néant). L'{\em an\={a}tman} est une
des caractéristiques de la souffrance (cf. {\em du\d{h}khasatya}). Le
Mah\={a}y\={a}na étend l'{\em an\={a}tman} du H\={\i}nay\={a}na à tous
les phénomènes ({\em dharm\={a}s}), et parle alors de vacuité
({\em \'{s}\={u}nyat\={a}}) d'existence inhérente. \\

{\dn a\7{n}\381wryogyAn} {\em anuttarayogay\={a}na} \ [pf. {\em an}
+ {\em uttara}  + {\em yoga} + {\em y\={a}na}] n. vaj. «~Véhicule
des yogas insurpassables\fg, désignation collective des trois
{\em y\={a}na} supérieurs du Vajray\={a}na:
{\em mah\={a}yogay\={a}na}, {\em anuyogay\={a}na} et
{\em atiyogay\={a}na}. \\

{\dn a\7{n}yogyAn} {\em anuyogay\={a}na} \ [{\em anu} + {\em yoga}
+ {\em y\={a}na}] n. vaj. «~Véhicule du yoga suivant\fg, cinquième des
six {\em y\={a}na} du Vajray\={a}na (second des trois
{\em y\={a}na} supérieurs). \\

{\dn aEBDm\0EpVk} {\em abhidharmapi\d{t}aka} [pf. {\em abhi} +
{\em dharma} + {\em pi\d{t}aka}] m. n. bd. «~Corbeille au sujet de la
doctrine du Bouddha\fg, troisième et dernière section
du canon bouddhique ({\em tripi\d{t}aka}), présentant la philosophie,
la psychologie et une glose des enseignements du Bouddha et de ses
principaux disciples (gloses du {\em s\={u}trapi\d{t}aka}), classés
systématiquement. Ce compendium sert de référence au
H\={\i}nay\={a}na et au Mah\={a}y\={a}na, et permet,
grâce aux différentes versions conservées, d'apprécier les différences
entre les sectes bouddhiques. On désigne couramment
l'{\em abhidharmapi\d{t}aka} par son contenu:
l'{\em abhidharma}. $\|$ pâli {\em abhidhammapitaka}. \\

{\dn aEBq\?k} {\em abhi\d{s}eka} m. vaj. «~Onction\fg, cérémonie
d'intronisation au cours de laquelle l'impétrant accepte la tutelle
d'un {\em guru}, seul garant d'une pratique correcte des
{\em y\={a}na} tantriques (cf. {\em tantra}). \\

{\dn ah\0^^F\qq{t}} {\em arhant} hin. «~Vénérable\fg, pratiquant parvenu au
plus haut degré de réalisation selon le H\={\i}nay\={a}na. Il a
vaincu les passions ({\em kle\'{s}a}) et les Dix Liens du
{\em sa\.{m}s\={a}ra}: l'illusion de la personnalité, le doute,
l'attachement aux rites et aux règles, le désir des sens, la rancune,
le désir de non corporéité ou de corporéité subtile, l'orgueil,
l'agitation et l'ignorance. Il est censé atteindre le
{\em nirv\={a}\d{n}a} à sa mort. À la différence du {\em bodhisattva}
du Mah\={a}y\={a}na qui a renoncé à son propre Éveil au profit
des êtres sensibles, l'{\em arhant} recherche l'Éveil pour
lui-même. Le Mah\={a}y\={a}na ne reconnaît pas au {\em
nirv\={a}\d{n}a} de l'{\em arhant} la perfection de celui qu'atteindra
le {\em bodhisattva}. En Chine l'{\em arhant\/} resta très populaire
bien que le bouddhisme y fut Mah\={a}y\={a}na --- sans doute parce que
cet idéal était proche de celui de l'anachorète taoïste. \\

{\dn aEv\38DwA} {\em avidy\={a}} [pf. {\em a} + {\em vidy\={a}}]
f. bd. «~L'ignorance\fg, première des causes de
la souffrance selon l'énumération traditionnelle faite par le
Bouddha (cf. {\em prat\={\i}tyasamutp\={a}da}). Selon le
{\em y\={a}na} elle prend une perspective un peu différente. De
fa\c{c}on générale elle est l'illusion fondamentale qui consiste à
séparer réellement {\em ceci\/} de {\em cela\/}: la dualité --- racine
de la souffrance (cf. {\em \={a}ryasatya}). \\

{\dn aAy\0mAg\0} {\em \={a}ryam\={a}rga} [{\em \={a}rya} + {\em
m\={a}rga}] m. bd. «~Le noble chemin\fg, dernière des Quatre
Nobles Vérités ({\em \={a}ryasatya}) du Bouddha ({\em buddha}). Elle
proclame l'existence d'une discipline et d'une doctrine qui mènent à
la cessation de la souffrance: le Noble Octuple Sentier ({\em
a\d{s}\d{t}\={a}\.{n}gam\={a}rga}). \\

{\dn aAy\0s(y} {\em \={a}ryasatya} [{\em \={a}rya} + {\em satya}] n.
bd. «~Noble vérité\fg, les Quatre Nobles Vérités proclamées par
le Bouddha ({\em buddha}) après son Éveil. Ce sont dans
l'ordre: la vérité de la souffrance ({\em du\d{h}kha}), la vérité
de l'origine de la souffrance ({\em samud\={a}ya}), la
vérité de la cessation de la souffrance ({\em nirodha}) et la
vérité du noble chemin ({\em \={a}ryam\={a}rga}) qui mène à la
cessation de la souffrance. Cet énoncé prend la forme d'un schéma
médical. D'après le Ny\={a}natiloka, le Bouddha dit: «~Mais qu'est-ce
donc, ô moines, que la Noble Vérité de la douleur? La naissance est
douleur, la maladie est douleur, le chagrin, la peine, la souffrance, la
tristesse, le désespoir sont douleur; la non-obtention de ce que l'on
désire est douleur; en un mot: les cinq constituants de l'existence
sont douleur. \\
Mais qu'est-ce donc, ô moines, que la Noble Vérité de l'apparition de
la douleur? C'est le désir, générateur de renaissance, avec son
cortège d'envies et de convoitises, s'assouvissant tantôt ici, tantôt
là, désir des sens, désir d'exister, désir d'autodestruction. \\
Mais qu'est-ce, ô moines, que la Noble Vérité de l'extinction de la
souffrance? C'est précisément l'extinction, l'abandon, le
renoncement, la délivrance, le détachement complets de ce désir. \\
Mais qu'est-ce, ô moines, que la Noble Vérité du Noble Sentier menant
à cette extinction de la souffrance? C'est le Noble Octuple Sentier, à
savoir: compréhension parfaite, pensée parfaite, parole parfaite,
action parfaite, moyens d'existence parfaits, effort parfait,
attention parfaite, concentration parfaite.\fg  \\

{\dn a\3A3wA\3BDwmAg\0} {\em a\d{s}\d{t}\={a}\.{n}gam\={a}rga} [{\em
a\d{s}\d{t}a} + {\em a\.{n}ga} + {\em m\={a}rga}] m. bd. «~Sentier
en huit étapes\fg, le Noble Octuple Sentier, explication de la
quatrième Noble Vérité ({\em \={a}ryam\={a}rga}) qui affirme
l'existence d'une discipline et d'une doctrine menant à la cessation
de la souffrance. Ce sentier est constitué par: 1. {\em la
compréhension parfaite} ou connaissance des Quatre Nobles Vérités
({\em \={a}ryasatya}) et de l'impersonnalité de l'existence ({\em
an\={a}tman}); 2. {\em la pensée parfaite} ou tolérance et
bienveillance envers tous les êtres vivants; 3. {\em la parole
parfaite} ou éviter le mensonge, le bavardage et la médisance; 4. {\em
l'action parfaite}  ({\em \'{s}\={\i}la}) ou éviter les actions
immorales; 5. {\em les moyens d'existence parfaits} ou éviter toute
profession pouvant nuire à d'autres êtres vivants (boucher, chasseur,
soldat, etc.); 6. {\em l'effort parfait} ou développer en soi ce qui
est sain et éviter ce qui est karmiquement malsain; 7. {\em
l'attention parfaite} ({\em sam\={a}dhi}) ou la méditation (cf. {\em
dhy\={a}na}). cf. {\em \={a}ryasatya}. \\

{\dn upyogyAn} {\em upayogay\={a}na} [pf. {\em upa} + {\em yoga} +
{\em y\={a}na}] n. vaj. «~Véhicule du yoga intermédiaire\fg, second des
six {\em y\={a}na} du  Vajray\={a}na (ou second des trois {\em
y\={a}na} inférieurs). \\

{\dn upAdAn} {\em up\={a}d\={a}na} [pf. {\em upa} +
{\em \={a}d\={a}}] bd. «~Saisie\fg, neuvième des douze causes de la
souffrance (cf. {\em prat\={\i}tyasamutp\={a}da}) $\mid$
bd. Attachement. D'après l'Abhidharmako\'{s}a, on peut
distinguer quatre sorte  d'{\em up\={a}d\={a}na}: l'attachement aux
sens, l'attachement aux idées, l'attachement aux rites et aux règles,
ainsi que  l'attachement à l'illusion de l'ego (cf.
{\em an\={a}tman} et {\em arhant}). $\|$ pâli {\em up\={a}d\={a}na} \\

{\dn kzZA} {\em karu\d{n}\={a}} f. mah. «~Compassion\fg, vertu
dominante du {\em bodhisattva}, qui le distingue de l'idéal du
H\={\i}nay\={a}na, l'{\em arhant}, qui met surtout l'accent sur la
sapience ({\em prajñ\={a}}) pour parvenir à l'Éveil. Dans le
Mah\={a}y\={a}na et surtout le Vajray\={a}na, la compassion et la
sapience sont dites co-émergentes de la réalisation de la vacuité
({\em \'{s}\={u}nyat\={a}}) des phénomènes ({\em dharm\={a}s}). Ainsi
compassion et sapience
sont liées et doivent se manifester ensemble dans l'action du {\em
bodhisattva}. La compassion de ce dernier peut prendre une forme
violente si la situation l'exige (cf. {\em bodhisattva}). En effet,
{\em karu\d{n}\={a}} implique un aspect actif, intrépide, en prise
directe avec la réalité, que ne recouvrent pas les traductions de
«~miséricorde~» (souvent liée à la rédemption chrétienne du pécheur)
et de «~pitié~» (qui est souvent la rencontre de la peur avec la
souffrance d'autrui; ou la «~douce pitié~» complaisante de Rousseau,
ou encore le mépris chez Nietzsche, ou bien la justice envers les
animaux et dont seraient exempts les juifs d'après Schopenhauer). Le
bouddhisme tibétain (cf.  Vajray\={a}na) insiste particulièrement sur
la compassion. \\

{\dn km\0} {\em karma} n. vaj. Dans le {\em ma\d{n}\d{d}ala} des Cinq
Bouddhas, dernière des Cinq Familles de Bouddha ({\em
buddhakula}). Son bouddha est Amoghasiddhi, son énergie est l'énergie
de l'action. Sa qualité confuse est la jalousie, l'envie. La sagesse
associée est «~la sagesse de l'accomplissement de l'action\fg, qui
per\c{c}oit que les situations ont les moyens de s'accomplir par
elles-mêmes, si bien qu'il est inutile de se précipiter et de
s'agiter. Son élément est le vent, sa couleur le vert. {\em karma} est
lié au nord, au crépuscule, à l'été. cf. {\em mah\={a}mudr\={a}}. \\

{\dn km\0\qq{n}} {\em karman} [{\em k\d{r}} {\bf I} + sf. {\em man}]
m. n. bd. «~Action, acte\fg, enchaînement causal inhérent à l'action,
à la fois responsable et produit du cycle des renaissances, ou {\em
sa\.{m}s\={a}ra} (de par son étymologie, «~réincarnation~» est un
terme inapproprié dans le contexte bouddhique car il sous-entend que
les renaissances ont toujours lieu {\em dans} un corps de chair et
donc qu'il y a une âme distincte du corps qui l'investit). Mais cette
causalité n'est pas déterminisme absolu: si les actes passés
déterminent la forme de la  renaissance et nous placent donc dans des
situations conséquentes, ce {\em karman} n'implique pas nos réponses à
ces situations --- l'Éveil est ainsi possible à tout instant. Il n'est
donc pas le destin des fatalistes (car il n'y a pas de dieu en soi) et
n'est pas plus imposé «~de l'extérieur~» (car il n'y a pas d'ego en
soi). Il est amoral; dire ainsi qu'il fait (ou qu'il est le fait) que
les mauvaises actions sont punies est une erreur: la cause contient en
germe la  conséquence, et n'est donc pas distincte d'elle, comme les
deux faces de la même médaille. Une image traditionnelle est: si la
flamme (l'esprit) d'une bougie (le corps) est soufflée et «~passe~» à
une autre bougie, on peut comparer ce souffle au {\em karman}.


Bien que le Bouddha ait mis
en garde contre une étude approfondie du {\em karman}, l'Abhidharma
(cf. {\em abhidharma}) distingue trois sources de {\em karman}, ou
Trois Portes: la parole, l'action (physique) et la pensée. Il est
important de noter
que l'attitude de l'esprit au moment de l'action est déterminante
du point de vue de la production de {\em karman} (ainsi tuer par
mégarde un homme au moyen d'une épée peut être moins «~négatif~» que de
transpercer un sac de grains en imaginant transpercer vraiment un
homme). De même un acte con\c{c}u dans l'esprit mais non concrétisé
aura des conséquences karmiques (en conséquence de l'intention).

On distingue un {\em karman} avec écoulement ({\em \={a}svara}) et un
sans écoulement. Le premier est créé par les passions ({\em
kle\'{s}a}) et enchaîne au {\em sa\.{m}s\={a}ra}. Selon la perspective
du H\={\i}nay\={a}na, le {\em karman} libérateur est celui sans
écoulement, et demande l'observation de préceptes moraux (ne pas tuer,
ne pas voler, ne pas commetre d'adultère, ne pas mentir, ne pas
calomnier). Le Mah\={a}y\={a}na défend le point de vue «~ni
écoulement, ni non-écoulement\fg, l'accumulation de mérites comme la
non-accumulation enchaînant à la dualité ({\em sa\.{m}s\={a}ra}/{\em
nirv\={a}\d{n}a}). À ce sujet, {\em L'extinction de la contemplation
[{\em vipa\'{s}yan\={a}}]}, attribué au {\em m\={a}dhyamika} chinois
Nieou-T'eou (583-657) dit (chap. X):

«~--- Existe-t-il des causes et des conditions permettant de tuer?

--- Le feu embrase les montagnes, le vent violent brise les arbres,
les avalanches ensevelissent les animaux sauvages, les eaux
débordantes noient les animaux rampants. Avec un esprit tel, l'on peut
tuer; mais si la moindre hésitation s'élève, avec la notion de vie et
de mort, il y a encore un esprit qui n'est pas épuisé, et même une
fourmi peut alors lier votre vie.

--- Existe-t-il des causes et des conditions permettant de voler et de
piller?

--- L'abeille butine la fleur de l'étang, le moineau picore la
châtaigne dans la cour, le buffle se nourrit de fèves, le cheval
broute l'herbe de la prairie. Sans l'idée de possession, même un pic
vous appartient; sinon la seule pointe d'une aiguille ou la moindre
petite feuille vous lie le cou et vous asservit.

--- Existe-t-il des causes et des conditions permettant la luxure?

--- Le ciel couvre la terre, le Yang s'unit au Yin, les latrines
re\c{c}oivent un écoulement du haut, les sources se déversent dans les
rigoles. Avec un esprit tel, tous les actes sont sans
obstruction. Mais si les passions engendrent la différenciation, alors
même votre propre femme souille votre esprit. [On ne peut s'empêcher
de penser au cynique Diogène à qui l'on reprocha de fréquenter les
bordels et qui répondit: «~Le soleil pénètre bien dans les latrines
sans se souiller!\fg]

--- Existe-t-il des causes et des conditions permettant le mensonge?

--- Si les paroles sont sans locuteur et les mots formulés sans donner
existence à l'esprit, la voix résonne comme la cloche, le souffle
bruit comme le vent. Avec un esprit tel, même ce que l'on nomme
Bouddha n'existe pas; sinon, même l'appellation de Bouddha est un
mensonge.\fg

Les {\em bodhisattva} doivent ainsi éliminer les notions de bon et
mauvais {\em karman}, puis éliminer la notion de «~ni bon ni mauvais
{\em karman}\fg.

Le terme {\em karman} est
traduit par «~karma~» (accepté par Le Petit Robert) lorsqu'il s'agit
d'insister sur l'aspect d'accomplissement de l'acte, sinon par
«~acte~» ou «~action\fg.\\

{\dn kSyAZ} {\em kaly\={a}\d{n}a} np. Ascète indien de Taxila,
dans la vallée de l'Indus, qui se joignit aux compagnons d'Alexandre
le Grand lors de leur retour en Grande Grèce. On pense qu'il
exer\c{c}a une influence sur Pyrrhon d'Élis, fondateur du scepticisme
antique. En Perse, se voyant faiblir, il pria Alexandre qu'on lui
élève un bûcher où il s'immola devant l'armée réunie. Dans les textes
grecs il est connu sous le nom de $K \acute{\alpha} \lambda \alpha \nu
o \varsigma$.\\

{\dn kSyAZEm/} {\em kaly\={a}\d{n}amitra} [{\em kaly\={a}\d{n}a}
+ {\em mitra}]  n. mah. «~Noble ami, ami dans le bien, ami
spirituel\fg, désigne le maître spirituel dans le Mah\={a}y\={a}na. Le
terme {\em guru} est employé lorsque la relation entre le pratiquant
et le maître devient très intense, voire despotique, comme dans le
Vajray\={a}na. \\

{\dn E\387wyAyogyAn} {\em kriy\={a}yogay\={a}na} [{\em kriy\={a}} +
{\em yoga} + {\em y\={a}na}] n. vaj. «~Véhicule du yoga de l'action\fg,
premier des {\em y\={a}na} du Vajray\={a}na (ou premier des
trois {\em y\={a}na} inférieurs). Il est lié au nettoyage des
dernières névroses, qui ne sont ni transmutables ni exploitables,
grâce aux {\em mantras}. syn. {\em kriy\={a}yoga}, {\em kriy\={a}yogatantra}. \\

{\dn "AE^^Ft} {\em k\d{s}\={a}nti} [{\em k\d{s}am}]
f. bd. «~Constance\fg, une des Six Perfections ({\em
p\={a}ramit\={a}}) du {\em bodhisattva}. On préfèrera «~constance~» à
«~patience\fg, car cette dernière peut être résignation due à une
frustation ou résolution liée à la volonté, ce qui est étranger à {\em
k\d{s}\={a}nti}. \\

{\dn \7{g}zyog} {\em guruyoga} [{\em gura} + {\em yoga}] m. vaj. «~Yoga
du gourou\fg, pratique tantrique ({\em yoga}) visant à développer la
dévotion fusionnelle avec son {\em guru}, comme moyen de progression
supérieur dans le Vajray\={a}na. \\

{\dn jAEt} {\em j\={a}ti} [vr. {\em jan} + sf. {\em ti}] f. bd. «~La
naissance\fg, onzième des douze causes de la souffrance. cf. {\em
prat\={\i}tyasamutp\={a}da}. \\

%{\dn \3E2wAn} {\em j\~{n}\={a}na} [{\em j\~{n}\={a}} + sf. {\em na}]
%n. bd. \\

{\dn tTAgt} {\em tath\={a}gata} [{\em tath\={a}} + {\em gata}]
m. bd. épith. «~Parvenu-au-oui~» (d'après les premières paroles du Bouddha
({\em buddha}) après l'Éveil: «~Oui, c'est ainsi.\fg), interprété
ultérieurement par l'Abhidharma ({\em abhidharma}) comme
«~allé-de-même~» (comme ses prédécesseurs) ou «~venu-de-même\fg. Dans
les {\em s\={u}tra}, le Bouddha est ainsi désigné ou se désigne
lui-même. \\

{\dn tTAgtgB\0} {\em tath\={a}gatagarbha} [{\em tath\={a}gata} +
{\em garbha}] m. mah. «~Germe du Parvenu-au-oui [le Bouddha],
Nature-de-Bouddha\fg. La nature éveillée fondamentale de tous les
êtres, temporairement voilée par les confusions dualistes. L'image
traditionnelle est celle du soleil derrière les nuages, ou d'un joyau
dans un tas de fumier. On utilise toujours la traduction moins
littérale de «~Nature-de-Bouddha\fg. cf. {\em tath\={a}gata} \\

%{\dn t^^F/} {\em tantra} [{\em tan} + sf. {\em tra}] \
%n. bd. «~Continuité\fg. \\

{\dn t^^F/yAn} {\em tantray\={a}na} \ [{\em tantra} + {\em
y\={a}na}] n. bd. Bouddhisme tantrique. Cette dénomination est utile
lorsque l'on veut distinguer le Vajray\={a}na au sein des {\em
y\={a}na}. syn. {\em vajray\={a}na}, {\em mantray\={a}na}.\\

{\dn \9{t}^^IZA} {\em t\d{r}\d{s}\d{n}\={a}} [{\em t\d{r}\d{s}} +
sf. {\em na}] f. bd. «~Le désir\fg, huitième des douze causes de la
souffrance. cf. {\em prat\={\i}tyasamutp\={a}da}. \\

{\dn E/kAyA\qq{s}} {\em trik\={a}y\={a}s} [pf. {\em tri} + {\em k\={a}ya}]
m. pl. mah. «~Les Trois Corps \fg, doctrine distinguant trois corps du
Bouddha mais qui ne font qu'un en réalité. Le premier est le {\em
dharmak\={a}ya}, corps du {\em dharma}, qui est la nature propre et
essentielle, la seule réalité, que l'école Madhyamaka considère comme
la vacuité [{\em \'{s}unyat\={a}}].  C'est le corps spirituel des
Bouddhas. Le second corps, {\em sa\d{m}bhogak\={a}ya}, est le corps de
félicité ou de fruition car il est acquis par les bonnes actions
antérieures. Le troisième corps, {\em nirm\={a}\d{n}ak\={a}ya}, est le
corps de transformation ou corps d'apparition dans lequel apparaît le
Bouddha pour convertir les êtres. Le zen explique les relations
qu'entretiennent ces trois corps à l'aide d'une analogie médicale: le
{\em dharmak\={a}ya} est comparable au savoir médical, le {\em
sa\d{m}bhogak\={a}ya} à la formation que le médecin re\c{c}oit en vue
d'acquérir ce savoir, et le {\em nirm\={a}\d{n}ak\={a}ya} à
l'application de ce savoir acquis dans le traitement des patients qui
recouvrent aisni la santé. \\

{\dn E/EpVk} {\em tripi\d{t}aka} [pf. {\em tri} + {\em
pi\d{t}aka}] m. n. bd. «~Les trois corbeilles\fg, le canon
bouddhique. Il se subdivise en trois parties: le {\em
vinayapi\d{t}aka} (discipline monastique), le {\em
s\={u}trapi\d{t}aka} (les dits du Bouddha) et l'{\em
abhidharmapi\d{t}aka} (les gloses des {\em s\={u}tra} et la
philosophie).\\

{\dn dAEknF} {\em d\={a}kin\={\i}} vaj. Déité féminine irritée ou
semi-irritée visualisée dans les pratiques ({\em s\={a}dhana}),
et parèdre des {\em heruka}. Elle représente la compassion ({\em
karu\d{n}\={a}}), la vacuité ({\em \'{s}\={u}nyat\={a}}) et la
sapience ({\em prajñ\={a}}). Les {\em d\={a}kin\={\i}} sont rusées et
joyeuses, elles représentent l'espace fondamental de la fertilité
({\em dharmadh\={a}tu}) d'où surgit le jeu du {\em sa\.{m}s\={a}ra} et
du {\em nirv\={a}\d{n}a}. Plus généralement, une {\em d\={a}kin\={\i}} peut
être une messagère ou une protectrice. \\

{\dn dAn} {\em d\={a}na} [{\em d\={a}} {\bf I} + sf. {\em na}]
n. bd. «~Générosité\fg, une des Six Perfections ({\em
p\={a}ramit\={a}}) que pratique le {\em bodhisattva}. Il est
recommendé aux {\em bodhisattva} dans leur pratique de {\em d\={a}na}
«~d'offrir leur don~» pour ne pas s'y attacher, et donc ne pas créer
de {\em karman}. C'est le sens des fréquentes dédicaces des mérites
après les méditations ({\em dhy\={a}na}) et à la fin des traités ({\em
\'{s}\={a}stra}). Ainsi {\em d\={a}na} n'est pas la charité chrétienne
qui elle constitue une bonne action: le {\em bodhisattva} lui ne recherche
ni le bien ni le mal et donc ne s'y attache pas; il agit comme la
situation l'exige, vue par la sapience ({\em praj\~{n}\={a}}). \\

{\dn \7{d},K} {\em du\d{h}kha} [pf. {\em dus} + {\em kha}]
n. bd. «~Souffrance\fg. Douleur physique et l'émotion pénible qui
l'accompagne. «~Malaise physique et mental, l'effort pour en sortir
et la dispersion qui en est la conséquence.~» (L. Silburn). $\mid$
n. bd. Souffrance fondamentale, inhérente à l'ego. Il
s'agit ici de la douleur de base, sous-jacente même aux sensations
plaisantes car celles-ci prendront toujours fin. Les dieux ({\em deva})
eux-mêmes n'en sont donc pas exempts car ils sont une forme de
renaissance, {\em donc} temporaire. cf. {\em du\d{h}khasatya}. \\

{\dn \7{d},Ks(y} {\em du\d{h}khasatya} [{\em du\d{h}kha} + {\em
satya}] n. bd. «~Vérité de la souffrance\fg, première des Quatre Nobles
Vérités ({\em \={a}ryasatya}). Elle consiste à constater que
l'existence est éphémère ({\em anitya}), impersonelle ({\em
an\={a}tman}) et, par conséquent, souffrance ({\em du\d{h}kha}). \\

{\dn Dm\0} {\em dharma} [{\em dh\d{r}} + sf. {\em ma}]
m. n. bd. «~Bonne Loi\fg, la voie bouddhique. D'aucuns traduisent {\em
dharma\/} par «~bouddhisme\fg, ce qui est {\em stricto sensu} incorrect
(cf. {\em triratna}). Il faut préférer «~bouddhisme~» pour se référer
à la philosophie, à la
psychologie, en bref au système (suivant en cela le suffixe -isme),
par opposition à {\em dharma} (non traduit) qui désigne alors la voie, la
pratique bouddhique proprement dite. $\mid$ bd. pl. {\em dharm\={a}s},
«~Phénomènes\fg, au sens de «~ce qu'il y a, ce qui est manifeste, qui est
per\c{c}u\fg. Bien que «~phénomène~» soit la traduction
universellemment adoptée, il faut l'entendre en fait au sens
d'«~apparence\fg, car le phénomène dans la philosophie occidentale
(phénoménisme d'Héraclite, scepticisme tardif, Nouvelle Académie,
Kant, phénoménologie, etc.)
laisse place à une éventuelle essence, à une substance, aux choses
(connaissables ou non, etc.). Pour les écoles du
H\={\i}nay\={a}na les agrégats ({\em skandha}) sont des phénomènes
(c'est-à-dire, en contexte bouddhique, dépourvus d'existence
intrinsèque --- cf. {\em an\={a}tman}), l'existence réelle des choses
n'étant pas explicitement niée (on parle alors parfois de
«~réalisme~» --- par excès car l'existence réelle des choses n'est
pas affirmée pour autant, et le principe du tiers exclus est dépourvu
de portée ontologique dans le bouddhisme.). Le Mah\={a}y\={a}na quant à
lui va plus loin et affirme que tout est phénomène et insiste
bien sur le fait que la nature des phénomènes est la vacuité
(d'existence intrinsèque) et que la vacuité est elle-même vide (Le
XIV$^e$ Dalaï Lama explique ce dernier point en disant que la vacuité
est nominale et que cette nominalité est elle-même nominale.). Ainsi
le Mah\={a}prajñ\={a}p\={a}ramit\={a}h\d{r}dayas\={u}tra («~Grand
soûtra du c{\oe}ur de la perfection de sapience\fg, vulgarisé sous le nom de
{\em Soûtra du c{\oe}ur}) dit: «~La forme [{\em dharma}] n'est que
le vide [{\em \'{s}unyat\={a}}], et le vide n'est rien d'autre que
la forme.~» Pour la traduction, il faut être conscient que du point
de vue philosophique occidental le bouddhisme ne peut être qualifié de
phénoménisme (\og apparence-{\em de}~» d'un être, de quelque chose)
ni de relativisme (\og apparence-{\em pour}~» un être, un sujet),
bien qu'il parle de phénomènes et de relativité. Ainsi, sur ce dernier
point, la co-production en dépendance (ou conditionnée) des phénomènes
est souvent qualifiée de \og relative~»pour signifier, par exemple,
que la cause ne se manifeste (et donc ne s'impute logiquement) qu'en
dépendance de la conséquence, et vice-versa --- les deux n'existant
pas séparément, donc intrinsèquement. N\={a}g\={a}rjuna, le fondateur
de l'école des Tenants du Milieu ({\em m\={a}dhyamika}), est parfois
qualifié de relativiste lorsqu'il rejette les paires d'opposés
(existence/inexistence par exemple) parce que chaque extrême n'existe
qu'en dépendance de l'autre, et donc lorsqu'il laisse non dit (il ne
peut vraiment \og affirmer~» puisque le principe de contradiction vient
de voler en éclats!) que la vérité inconditionnée ou ultime ({\em
param\={a}rthasatya}) des phénomènes réside entre les deux extrêmes,
au milieu ({\em m\={a}dhyama}), c'est-à-dire dans la vacuité ({\em
\'{s}\={u}nyat\={a}}). Dans aucun de ces deux cas de figure il ne
s'agit de relativisme au sens où l'entend la philosophie
occidentale. $\|$ p\={a}li {\em dhamma}. \\

{\dn Dm\0DA\7{t}} {\em dharmadh\={a}tu} [{\em dharma} + {\em
dh\={a}tu}] m. mah. \og l'Espace-qui-englobe-tout\fg, la toile de fond
de la réalité\fg, le Domaine absolu, l'Espace foncier, l'Espace
fondamental de la fertilité~» d'où émerge le jeu du {\em
sa\.{m}s\={a}ra} et du {\em nirv\={a}\d{n}a}. \\

{\dn @yAn} {\em dhy\={a}na} [{\em dhy\={a}} + sf. {\em na}]
n. mah. \og Méditation\fg, quatrième des Cinq Catégories ({\em bala}) de
la voie de l'unification. $\mid$
n. bd. Méditation, absorption
dépourvue de toute saisie (mentale), ou bien concentration
analytique. Bien qu'universellement adoptée, la traduction
fran\c{c}aise \og méditation\fg\
risque le contresens car le mot a le sens de \og profonde réflexion\fg\
(cf. Descartes, Lamartine, etc.), ce qui peut être très différent de
{\em dhy\={a}na} (notamment s'il s'agit de {\em \'{s}amatha}, dans le
bouddhsime tibétain, ou de {\em zazen} dans le bouddhisme
Zen S\={o}t\={o}). Il n'est pas besoin de lieu particulier pour
pratiquer; ainsi D\={o}gen Zenji (maître zen) écrivit: \og Ton propre
c{\oe}ur, voilà ta salle d'exercice.\fg (le {\em c{\oe}ur} désigne
dans le bouddhisme le siège de l'esprit --- au sens du {\em pneuma}
grec, et non du {\em mens} latin --- et pas celui des
sentiments). Contrairement à ce que l'on croit fréquemment, {\em
dhy\={a}na} n'implique pas forcément une  assise immobile: elle peut
se  pratiquer pendant les quatre activités du corps (selon un texte du
Mah\={a}y\={a}na): assis, marchant,  couché et debout. Les Patriarches
Ch'an par exemple ne donnaient pas de consignes particulières sur la
posture à adopter.
Aujourd'hui toutes les traditions bouddhiques réclament une assise
immobile, mais comme point d'ancrage de la pratique et jamais comme
une fin en soi: il s'agira toujours ensuite de \og mettre la
méditation en action~» pour qu'elle ne fasse plus qu'un avec tous
les aspects de la vie phénoménale (courante). La méditation est avant
tout une attitude de l'esprit {\em et} du corps sur le chemin de
l'Éveil, ou qui est l'expression même de l'Éveil (les bouddhistes zen
préfèrent répéter inlassablement que {\em dhy\={a}na} ---  {\em zazen}
--- {\em est}
l'Éveil, qu'il ne s'agit donc aucunement d'un moyen; ils sont fidèles
en cela à leur tradition subitiste). En Chine, Houai-jang (677-744)
dit à son disciple Ma-tsou: \og Désires-tu apprendre à être assis en
{\em dhy\={a}na}, ou bien à être assis en Bouddha? Si tu veux apprendre
à être assis en {\em dhy\={a}na}, sache que {\em dhy\={a}na} ne relève
ni de la position assise ni de la position allongée. Si tu veux
apprendre à être assis en Bouddha, sache que le Bouddha n'a pas de
caractéristiques déterminées. Au sein du {\em dharma} de la
non-demeure, ne t'occupe ni d'obtention ni d'attachement. Si tu
t'assois en Bouddha, tu tues le Bouddha. Si tu t'attaches à la
position assise, tu n'atteindra pas la vérité absolue.~»Mais on
retiendra aussi que Milarepa en guise d'ultime enseignement à son
disciple Gampopa lui montra ses fesses couvertes de cals, fruits d'une
pratique assidue de {\em dhy\={a}na}. cf. {\em \'{s}amatha}, {\em
vipa\'{s}yan\={a}}, {\em sam\={a}dhi}. \\

{\dn mnomykAy} {\em manomayak\={a}ya} [{\em man} + {\em maya} +
{\em k\={a}ya}]. m. mah. \og Corps [fait] de pensée\fg, corps
produit pour sauver les êtres. Les {\em bodhisattva} de la dixième
Terre ({\em bh\={u}mi}) en ont la capacité. Il est décrit dans le
La\.{n}k\={a}vat\={a}ras\={u}ra: \og Le Bienheureux dit alors: `Il y a
trois sortes de corps produits par la pensée, Mah\={a}mati. Quels sont
ces trois? Il y a le corps obtenu dans la jouissance du {\em
sam\={a}dhi}, le corps obtenu par la reconnaissance intrinsèque de
toute chose, et le corps assumé par un {\em bodhisattva} selon la
catégorie d'êtres qu'il a à sauver. [...] En ce qui concerne toutes
les sortes de corps produits par la pensée, je dis qu'ils ne sont rien
d'autre que des évaluations de l'esprit'.\fg \ On retrouve là exactement
le {\em trik\={a}ya}. \\

{\dn nAm!p} {\em n\={a}mar\={u}pa} [{\em n\={a}ma} + {\em
r\={u}pa}] n. bd. \og Le nom et la forme\fg, quatrième des douze causes
de la souffrance. cf. {\em prat\={\i}tyasamutp\={a}da}. \\

%{\dn EndAn} {\em nid\={a}na} bd. \\

%{\dn EnvA\0Z} {\em nirv\={a}\d{n}a} [{\em nirv\={a}} + sf. {\em na}]
%n. bd. $\|$ p\={a}li {\em nibb\={a}na}. \\

{\dn p\394w} {\em padma} m. n. bd. \og Lotus\fg, symbole
du chemin vers l'Éveil, car il pousse dans les eaux boueuses et éclot
à la lumière du jour dans une blancheur immaculée. Dans
l'iconographie le bouddha \'{S}\={a}kyamuni, ainsi que certains
bouddhas et bodhisattvas transcendantaux, est représenté assis dans
une fleur de lotus (si elle possède huit pétales elle évoque alors le
Noble Octuple Sentier ou {\em a\d{s}\d{t}\={a}\.{n}gam\={a}rga}). $\mid$
m. n. vaj. Dans le {\em ma\d{n}\d{d}ala} des Cinq Bouddhas, quatrième
des Cinq Familles de Bouddha ({\em buddhakula}). Son bouddha est Amit\={a}bha,
son énergie est le magnétisme. Sa qualité confuse est la passion. La
sagesse qui lui est associée est \og la sagesse discriminante\fg, qui
consiste à  voir les détails des situations avec une infinie
précision. Son élément est le feu, sa couleur le rouge. {\em padma}
est lié à l'ouest, au printemps. cf. {\em mah\={a}mudr\={a}} \\

{\dn prmAT\0s(y} {\em param\={a}rthasatya} [{\em parama} +
{\em artha} + {\em satya}] n. mah. \og Vérité ultime de la réalité\fg.
Vérité (ou sens) ultime, une des deux catégories que distingue la
logique du Mah\={a}y\={a}na pour classer les objets de
connaissance (l'autre est la vérité conventionnelle, ou
{\em sa\d{m}v\d{r}tisatya}). \og Une vérité conventionnelle est un objet
per\c{c}u de fa\c{c}on dualiste par une connaissance valide directe
qui le per\c{c}oit en fait. L'absence d'existence réelle d'un vase est
un exemple de vérité ultime, le vase un exemple de vérité
conventionnelle. Ceci est valable pour tous les phénomènes [{\em
dharm\={a}s}], tous sans exception possèdent ces deux vérités. [...]
Une vérité ultime est une vacuité [{\em \'{s}\={u}nyat\={a}}]. [...]
Si on les divise, les
vérités ultimes sont au nombre de seize ou, en bref, de quatre
vacuités. Ces classifications sont établies par rapport aux bases sur
lesquelles s'exercent la conception d'un soi, d'une existence réelle:
intérieures, extérieures, la vacuité elle-même, etc. Elles servent à
contrecarrer cette conception mais n'impliquent aucune différence
quant à la vacuité proprement dite.~» (G. Khensur) \\

{\dn pErEnvA\0Z} {\em parinirv\={a}\d{n}a} n. bd. \og Extinction
parfaite\fg. C'est le {\em nirv\={a}\d{n}a} complet et définitif qui
n'est atteint qu'à la mort. C'est pourquoi dans les textes ce terme
est synonyme de la mort du bouddha \'{S}\={a}kyamuni. \\

{\dn pArEmtA} {\em p\={a}ramit\={a}} \ [{\em p\={a}ra} +
sf. {\em it\={a}}]  mah. \og Rejoignant l'autre rive\fg, les Six
Perfections que pratique le {\em bodhisattva\/}: générosité
({\em d\={a}na}), discipline ({\em \'{s}\={\i}la}), constance ({\em
k\d{s}\={a}nti}), énergie ou vigueur ({\em v\={\i}rya}), méditation
({\em dhy\={a}na}) et sapience ({\em prajñ\={a}}). \\

{\dn \3FEw\3E2wA} {\em praj\~{n}\={a}} \ [{\em pra} + {\em jñ\={a}}]
f. mah. \og Sapience, sagesse, connaissance transcendante\fg, dernière des
Cinq Catégories ({\em bala}) de la voie de l'unification. Il s'agit d'une
connaissance pénétrante  de la réalité. Le bouddhisme,
contrairement aux philosophies occidentales, n'envisage pas la
connaissance mais {\em l'acte de prendre connaissance}. Il faut donc
entendre ici: \og Prendre/Faire connaissance de fa\c{c}on pénétrante
de/avec la réalité\fg. Dans la perspective tantrique (cf. {\em
vajaray\={a}na}), la {\em prajñ\={a}} est la mère de tous les bouddhas,
celle qui donne naissance à l'idée même d'Éveil. Selon le
Mah\={a}y\={a}na elle est réalisée à la sixième Terre ({\em bh\={u}mi}), ou
palier, de la voie du {\em bodhisattva}. La logique bouddhique
est très différente de la logique aristotélicienne en ce sens qu'elle ne
cherche pas à saisir une vérité (en elle-même) mais, par le biais d'une
dialectique d'inférences négatives, à détourner l'esprit des vues
erronées --- le langage, ou discours ({\em logos} en grec), ne pouvant
saisir ni contenir la réalité ultime. Il semble que ce soit là un
point commun avec le scepticisme tardif occidental (Sextus
Empiricus). Certains auteurs évitent la traduction
\og sagesse~» à cause de la connotation profane, et donc presque
péjorative, que la philosophie occidentale lui a donné. La
transcendance {\em en opposition inconciliable avec l'immanence} étant
étrangère au bouddhisme (elle a été de plus très connotée par les
philosophes occidentaux), sans doute faut-il préférer \og sapience\fg,
terme de la gnose qui comprend les deux aspects de connaissance et de
sagesse. Le terme de \og sagesse~»sera alors réservé à {\em
jñ\={a}na}. $\mid$ f. mah. La sapience, dernière des Six Perfections
({\em p\={a}ramit\={a}}) du {\em bodhisattva}. Elle est l'élément de
clarté et de précision présent dans les autres perfections. \\

{\dn \3FEwtF(ys\7{m}(pAd} {\em prat\={\i}tyasamutp\={a}da} \
bd. Production en chaîne de la souffrance ({\em du\d{h}kha}), qui
soutend la seconde des Nobles Vérités ({\em samud\={a}ya}). Cette
chaîne fermée est formée de douze maillons ou liens ou causes ou
conditions de production ({\em nid\={a}na}), et peut être parcourue
dans les deux sens: ignorance ({\em avidy\={a}}), action volitive
({\em sa\.{m}sk\={a}ra}), conscience ({\em vij\~{n}\={a}na}), nom et
forme ({\em n\={a}mar\={u}pa}), perception sensorielle ({\em
\d{s}a\d{d}\={a}yatana}), toucher ({\em spar\'{s}a}), sensation ({\em
vedan\={a}}), désir ({\em t\d{r}\d{s}\d{n}\={a}}), saisie ({\em
up\={a}d\={a}na}), devenir ({\em bh\={a}va} --- caractérisé par le
coït, une autre saisie), naissance ({\em j\={a}ti}), puis vieillesse
et mort ({\em jar\={a}}). Traditionnellement leur symboles sont
respectivement: une grand-mère aveugle, un tour de potier, un singe,
une personne dans un bateau, un singe dans une maison à six fenêtres,
un couple marié, une flêche dans l'{\oe}il, boire du lait et du miel,
cueillir des fruits, coïter, une femme en train d'enfanter, un cortège
funèbre. cf. {\em \={a}ryasatya}. \\

%{\dn \3FEw(y\?k\7{b}\388w} {\em pratyekabuddha} m. bd. \\

{\dn \3FEw(y\?k\7{b}\388wyAn} {\em pratyekabuddhay\={a}na} [{\em
pratyeka} + {\em buddha} + {\em y\={a}na}] n. mah. \og Véhicule du
Bouddha qui s'éveille de lui-même\fg, dans la perspective du
Mah\={a}y\={a}na, dernier des deux {\em y\={a}na} du
H\={\i}nay\={a}na. L'accent est mis sur sur le travail sur soi. C'est
la voie de la compréhension psychologique (donc interne, par rapport
au {\em \'{s}ravakabuddhay\={a}na}) de la nature de la vie, des Cinq
Agrégats ({\em skandha}). \\

{\dn bl} {\em bala} n. mah. \og Force\fg, les Cinq Pouvoirs ou
Catégories. Le développement de ces cinq principes caractérise la voie
de l'unification: la confiance inébranlable ({\em \'{s}raddh\={a}}),
l'énergie ou vigueur ({\em v\={\i}rya}), l'attention ou recueillement
({\em sm\d{r}ti}), la méditation ({\em sam\={a}dhi}) et la sapience
({\em prajñ\={a}}). \\

{\dn bFjm^^F/} {\em b\={\i}jamantra} \ [{\em b\={\i}ja} +
{\em mantra}] m. n. vaj. \og Syllabe-germe\fg, une syllabe, généralement
en sanscrit, qui représente la réalité essentielle d'une déité
particulière visualisée par le pratiquant tantrique (pratique nommée
{\em s\={a}dhana}). Les tibétains employent des syllabes
tibétaines. cf. {\em vajray\={a}na}, {\em d\={a}kin\={\i}}. \\

{\dn \7{b}\388w} {\em buddha} \ bd. épith. du bouddha \'{S}\={a}kyamuni
\og l'Éveillé\fg; né prince Siddh\={a}rtha Gautama à Kapilavastu du roi
{\em \'{s}\={a}kya} \'{S}uddhodana et de la reine M\={a}y\={a}
(env. 560 - env. 480 av. JC); à l'âge de 16 ans il épousa
Ya\'{s}odhar\={a} dont il eut un fils, R\={a}hula; à l'âge de 29 ans
il quitta son palais pour pratiquer l'ascétisme pendant 6 ans. Il
atteignit l'Éveil [{\em bodhi}] à Gay\={a} et proclama les Quatre
Nobles Vérités ({\em \={a}ryasatya}) dans le Parc aux gazelles de
Bénarès ({\em m\d{r}gad\={a}va}): il exécuta ainsi \og le premier tour
de la Roue de la Loi~» ({\em dharmacakra}), et fonda une communauté
monastique (la {\em sa\d{m}gha}). Il entra dans le  {\em
parinirv\={a}\d{n}a}, l'extinction parfaite, à l'âge de 80 ans. Il n'a
laissé aucun écrit. $\mid$ m. bd. Personne ayant réalisé l'Éveil ({\em
nirv\={a}\d{n}a}). $\mid$ vaj. Dans le {\em ma\d{n}\d{d}ala} des Cinq
Bouddhas, première des Cinq Familles de Bouddha ({\em buddhakula}),
usuellement  placée au centre du {\em ma\d{n}\d{d}ala}. Son bouddha
est Vairocana, son énergie est la contemplation. Sa qualité confuse
est  la stupidité paresseuse. La sagesse qui lui est associée est \og
la sagesse du  {\em dharmadh\={a}tu}~» qui pénètre et active les
quatre autres sagesses. {\em buddha} est lié à la couleur bleue, à la
qualité fraîche et vaste du ciel. cf. {\em mah\={a}mudr\={a}}.\\

{\dn \7{b}\388w\7{k}l} {\em buddhakula} [{\em buddha} + {\em kula}]
n. vaj. \og Familles de bouddha\fg, les cinq qualités du {\em
sa\d{m}bhogak\={a}ya} telles qu'elles apparaissent dans le {\em
ma\d{n}\d{d}ala} des Cinq Bouddhas. cf. {\em
mah\={a}mudr\={a}}. \\

{\dn boEDEc\381w} {\em bodhicitta} [{\em bodhi} + {\em citta}]
n. mah. \og Esprit d'Éveil\fg, \og Vaste intention de l'esprit d'Éveil\fg,
forte aspiration à atteindre l'Éveil pour venir en aide aux êtres
sensibles. Elle est le préalable à l'engagement sur le sentier du {\em
bodhisattva}. \\

{\dn boEDs\382wv} {\em bodhisattva} [{\em bodhi} + {\em sattva}]
m. mah. \og Être [courageux] pour l'Éveil, héros pour l'Éveil\fg, pratiquant
idéal du Mah\={a}y\={a}na qui, s'appliquant aux Six Perfections
({\em p\={a}ramit\={a}}), renonce à atteindre l'Éveil pour
lui-même et se consacre à aider les êtres doués de sensibilité sur le
chemin de l'Éveil. Par rapport à l'{\em arhant} du
H\={\i}nay\={a}na qui recherche l'Éveil pour lui-même, l'accent
est mis sur la compassion ({\em karu\d{n}\={a}}), et son action est
considérée comme héroïque et même parfois guerrière. Le {\em bodhisattva} ne
cherche pas à \og faire le bien\fg, en suivant une certaine morale, mais
à accomplir ce qui est juste, c'est-à-dire conforme à la situation,
telel que la per\c{c}oit la sapience ({\em praj\~{n}\={a}}). Sa
compassion peut ainsi s'exprimer de fa\c{c}on violente et impitoyable
s'il le faut, surtout dans la tradition tantrique (Vajray\={a}na) où il est
dit que \og si vous ne détruisez pas quand il le faut, vous manquez au
v{\oe}u de compassion qui vous engage formellement à détruire ce qui
est frivole. En conséquence, suivre la voie ne veut pas nécessairement
dire essayer seulement d'être bon et de n'offenser personne. Cela ne
fonctionne pas. Cela n'est pas la question. Si quelqu'un traverse
brusquement notre route, il faut le repousser parce que son intrusion
est sans objet.~» (C. Trungpa) Pour la traduction on préfèrera conserver
\og bodhisattva~» (accepté par Le Petit Robert). \\

%{\dn m^^WXl} {\em ma\d{n}\d{d}ala} m. n. bd. \\

%{\dn m^^F/} {\em mantra} [{\em man} + sf. {\em tra}]
%m. n. vaj. \og Protection de l'esprit\fg. \\

{\dn m^^F/yAn} {\em mantray\={a}na} [{\em mantra} + {\em y\={a}na}]
n. vaj. Bouddhisme tantrique. syn. {\em vajray\={a}na},
{\em tantray\={a}na}. \\

{\dn mhA\7{m}\qb{d}A} {\em mah\={a}mudr\={a}} [{\em mah\={a}} + {\em
mudr\={a}}] f. vaj. \og Grand Sceau\fg, un des enseignements du
Vajray\={a}na. {\em mah\={a}} doit être ici entendu dans le sens
d'\og insurpassable~» et {\em mudr\={a}} non pas représentation mais
actualisation, manifestation, des qualités vivantes. Le {\em
mah\={a}mudr\={a}} est l'enseignement du {\em ma\d{n}\d{d}ala} des
Cinq Bouddhas, où aux quatre points cardinaux et au centre
se trouvent placés cinq bouddhas du {\em sa\d{m}bhogak\={a}ya}, dits
\og Familles de bouddha~» ({\em buddhakula}): {\em buddha}, {\em vajra},
{\em ratna}, {\em padma}, et {\em karma}. Chacune représente un aspect
de la condition éveillée de l'esprit et son aspect confus (une
émotion, une qualité du {\em sa\.{m}s\={a}ra}). Le pratiquant doit
transmuter ces aspects confus en leur pendants éveillés. Des couleurs,
des éléments, des paysages, des saisons, tous les aspects du monde
phénoménal sont associés aux Familles de Bouddha.  cf. {\em
y\={a}na}. \\

%{\dn mhAyAn} {\em mah\={a}y\={a}na} n. bd. \\

{\dn mhAyogyAn} {\em mah\={a}yogay\={a}na} \ [{\em mah\={a}} +
{\em yoga} + {\em y\={a}na}]  n. vaj. \og Véhicule du yoga
insurpassable\fg, quatrième des six  {\em y\={a}na} du Vajray\={a}na
(ou premier des trois {\em y\={a}na} supérieurs). syn. {\em mah\={a}yoga},
{\em mah\={a}yogatantra}. \\

{\dn mhAEs\388w} {\em mah\={a}siddha} m. vaj. \og Grand accompli\fg. Se
dit de certains maîtres du bouddhisme tantrique, détenteurs de
pouvoirs magiques. \\

%{\dn mA@ymk} {\em m\={a}dhyamaka} [{\em madhyama} + sf. {\em
%aka}] mah. \og Voie du milieu\fg\\

%{\dn mA@yEmk} {\em m\={a}dhyamika} mah. \og Tenant de la Voie du
%Milieu\fg. \\

{\dn yAn} {\em y\={a}na} n. bd. \og Véhicule\fg. Ensemble cohérent
d'enseignements intellectuels et de méthodes pratiques de méditation
liés à un stade particulier de la progression du pratiquant sur le
chemin du Bouddha. Les trois principaux véhicules sont le
H\={\i}nay\={a}na, le Mah\={a}y\={a}na et le
Vajray\={a}na. Le H\={\i}nay\={a}na peut être subdivisé en deux
{\em y\={a}na}: le {\em \'{s}ravakay\={a}na} puis le
{\em pratyekabuddhay\={a}na}. Pour l'École de l'Ancienne
Traduction que suit l'ordre tibétain des Nyingma, le
Vajray\={a}na se subdivise en six {\em y\={a}na}: trois {\em
y\={a}na} inférieurs: {\em kriy\={a}yogay\={a}na}, {\em
upayogay\={a}na},  {\em yogay\={a}na}, et trois {\em
y\={a}na} supérieurs (désignés collectivement par {\em
anuttarayogay\={a}na}): {\em mah\={a}yogay\={a}na}, {\em
anuyogay\={a}na} et {\em atiyogay\={a}na}. Selon l'École de
la Nouvelle Traduction, que suivent les ordres tibétains Kagyü,
Sakya et Gelug, le Vajray\={a}na se subdivise en quatre {\em
y\={a}na}: {\em kriy\={a}yogay\={a}na}, {\em upayogay\={a}na}, {\em
yogay\={a}na},  {\em anuttarayogay\={a}na} (différent du {\em
anuttarayogay\={a}na} de l'Ancienne École). L'enseignement du {\em
mah\={a}mudr\={a}} appartient à cette dernière école. \\

%{\dn yogAcAr} {\em yog\={a}c\={a}ra} mah. \og L'exercice du yoga\fg.\\

{\dn yogyAn} {\em yogay\={a}na} \ [{\em yoga} + {\em y\={a}na}]
n. vaj. \og Véhicule des yogas\fg, troisième des six {\em y\={a}na} du
Vajray\={a}na (ou dernier des trois {\em y\={a}na} inférieurs). \\

{\dn r\3D7w} {\em ratna} [{\em r\={a}} {\bf I} + sf. {\em na}]
n. vaj. \og Joyau\fg. Dans le {\em ma\d{n}\d{d}ala} des Cinq {\em
tath\={a}gata}, troisième des Cinq Familles de Bouddha ({\em
buddhakula}). Son bouddha est Ratnasa\d{m}bhava, son énergie associée
est la richesse. Sa qualité confuse est l'orgueil. La sagesse qui lui
est associée est \og la sagesse de l'équanimité~» qui consiste à
embrasser toutes les situations d'un seul coup avec une vision
panoramique. Son élément est la terre, sa couleur le jaune. {\em
ratna} est lié au sud, à l'automne. cf. {\em mah\={a}mudr\={a}}. \\

{\dn !p} {\em r\={u}pa} n. bd. \og La forme\fg, premier des Cinq
Agrégats  ({\em skandha}): la forme, l'être fondamental qui est la
base de la dualité, de la distinction entre sujet et objet, projecteur
et projection. C'est l'ignorance fondamentale de notre
nature-de-bouddha  ({\em tath\={a}gatagarbha}), une ignorance active
et sans effort (comme l'instinct de saisie qui permet à l'oiseau de
dormir sur la branche) qui fait tout pour s'ignorer et est la
source des innombrables stratégies de l'ego en vue de sa
perpétuation. (C. Trungpa) \\

%{\dn lok} {\em loka} m. bd. \og monde\fg. \\

{\dn v\385w} {\em vajra} [{\em vaj} + sf. {\em ra}] m. n. vaj. Le
\og diamant\fg, en tant qu'objet précieux et indestructible, dont le
tranchant et la dureté viennent à bout des illusions. $\mid$ vaj. Dans
le {\em ma\d{n}\d{d}ala} des Cinq Bouddhas, seconde des Cinq Familles
de Bouddha ({\em buddhakula}). Son bouddha est Ak\'{s}obhya, son
énergie est l'intelligence. Sa qualité confuse est la colère. La
sagesse associée est \og la sagesse semblable au miroir~» qui est la
clarté où l'agression est vue comme un intellect précis,
tranchant. L'élément associé est l'eau, sa couleur le  blanc. {\em
vajra} est lié à l'est, à l'aube, à l'hiver.  cf. {\em
mah\={a}mudr\={a}}. \\

{\dn v\385wmAlA} {\em vajram\={a}l\={a}} [{\em vajra} + {\em
m\={a}l\={a}}] np. vaj. \og La guirlande de diamant\fg, texte du {\em
kriy\={a}yogay\={a}na} mettant notamment en garde contre les vues
erronées des visualisations de déités ({\em s\={a}dhana}). \\

{\dn v\385wyAn} {\em vajray\={a}na} [{\em vajra} + {\em y\={a}na}]
n. bd. \og Véhicule adamantin\fg, bouddhisme tantrique. Pour sa
légitimation doctrinale il se considère comme le troisième tour de la
Roue de la Loi ({\em dharmacakra}) donné par le Bouddha sous sa forme
de  {\em sa\d{m}bhogak\={a}ya}, et donc le prolongement naturel du
Mah\={a}y\={a}na. Il est réputé constituer une voie rapide, mais très
dangeureuse, d'accès à l'Éveil (en une vie), bien que, présupposant la
pratique du H\={\i}nay\={a}na et du Mah\={a}y\={a}na, il
insiste sur la nécessité d'une méthode progressive (sur ce
dernier point il s'oppose aux écoles Mah\={a}y\={a}na subitistes
du Ch'an et du Zen). Du point de vue scripturaire, le Vajray\={a}na se
caractérise par les {\em tantra} et l'utilisation de formules orales
dites {\em mantra}. Le Vajray\={a}na se concentre particulièrement sur
la transmutation des énergies vitales, des émotions et des névroses
en aspects de l'esprit éveillé (cf. {\em mah\={a}mudr\={a}}). L'accent
est mis aussi sur l'établissement d'une relation totale avec un maître
pleinement réalisé ({\em guru}) qui, par rapport au
Mah\={a}y\={a}na où il constitue un ami spirituel ({\em
kaly\={a}namitra}), est ici \og un despote éclairé~» (C. Trungpa) qui
exige confiance et abandon dans le contexte des liens créés par un
v{\oe}u appelé {\em samaya}. Le Vajray\={a}na pousse à son point ultime
la logique du Mah\={a}y\={a}na en reconnaissant pleinement
l'interdépendance du {\em nirv\={a}\d{n}a} et du {\em
sa\.{m}s\={a}ra}, et donc ne recherchant pas le premier en rejetant le
second: tous les aspects de la vie phénoménale sont des occasions
d'Éveil --- la sexualité, l'alcool, les excréments y compris. Son
apparente immoralité (il s'agit en fait d'amoralité), son ésotérisme
(il s'agit en fait de protéger le pratiquant néophyte et non d'exprimer une
domination), ses pratiques psycho-physiques (comme la visualisation de
déités, dite {\em s\={a}dhana}), son iconographie sexuelle et mortuaire,
son long isolement géographique (au Tibet surtout) et la
confusion avec le {\em tantra} hindouiste en font un bouddhisme
grandement incompris. syn. {\em tantray\={a}na}, {\em
mantray\={a}na}. \\

{\dn v\385ws\382wv} {\em vajrasattva} [{\em vajra} + {\em
sattva}] np. vaj. \og Héros de diamant\fg, l'une des déités visualisées
à différents niveaux de pratique (cf. {\em y\={a}na}). Elle est associée
à la pureté primordiale. \\

{\dn Ev\3E2wAn} {\em vij\~{n}\={a}na} [{\em vijñ\={a}} {\bf II} +
sf. {\em na}] n. bd. \og La conscience\fg, dernier des Cinq Agrégats
({\em skandha}): elle contient les émotions et les
modèles de pensée de toutes sortes. Les émotions proviennent de
la frustation, au sens où nous pourrions être incapables d'assouvir
nos désirs, ou le sommes déjà. Le méchanisme des émotions
double ainsi celui des projections (qui sont notre version des
phénomènes ou {\em dharm\={a}s}), qui apparaît alors suspect ou
ennuyeux. (C. Trungpa) $\mid$ bd. Troisième des douze causes de
la souffrance. cf. {\em prat\={\i}tyasamutp\={a}da}. \\

{\dn Ev\38DwADr} {\em vidy\={a}dhara} [{\em vidy\={a}} + {\em dhara}]
vaj. \og Détenteur de la folle sagesse\fg. Maître du bouddhisme tantrique. \\

{\dn EvnyEpVk} {\em vinayapi\d{t}aka} [{\em vinaya} + {\em
pi\d{t}aka}] m. n. bd. \og Corbeille de la discipline\fg, section du
canon bouddhique ({\em tripi\d{t}aka}) recueillant les prescriptions
du Bouddha ({\em buddha}) concernant la discipline monastique. \\

{\dn Evp[ynA} {\em vipa\'{s}yan\={a}} f. hin. Méditation ({\em
dhy\={a}na}) destinée à développer l'attention (cf. {\em
sm\d{r}tiupasth\={a}na}), et utilisée en alternance avec {\em
\'{s}amatha}. cf. {\em sam\={a}dhi}. $\|$ p\={a}li {\em
vipassan\={a}}. \\

{\dn vFy\0} {\em v\={\i}rya} [{\em v\={\i}ra} + sf. {\em ya}]
n. mah. \og La vigueur, l'énergie\fg, seconde des Cinq Catégories ({\em
bala}) de la voie de l'unification. Vigueur, énergie enjouée à
poursuivre la voie, qui naît de l'intuition que \og quelque chose se
passe\fg. \'{S}\={a}ntideva la définit comme \og Prendre plaisir au
mérite de se redécouvrir\fg. $\mid$ mah. La quatrième des Six
Perfections ({\em p\={a}ramit\={a}}) du {\em bodhisattva}. \\

{\dn v\?dnA} {\em vedan\={a}} f. bd. \og La sensation\fg, second des Cinq
Agrégats ({\em skandha}). Elle est prise de conscience de soi, à un
niveau élémentaire, par rapport à un environnement dont il s'agit
alors de déterminer, de contr\^oler, s'il nous est agréable,
désagréable ou indifférent (projections de base). (C. Trungpa) \\

{\dn fAE^^Ftd\?v} {\em \'{s}\={a}ntideva}  np. mah. (VII$^e$-VIII$^e$
siècle) Prince puis moine bouddhiste Mah\={a}y\={a}na et maître de
l'école des Tenants du Milieu {\em m\={a}dhyamaka}. Il est l'auteur du
Bodhisattvacary\={a}vat\={a}ra. \\

{\dn fmT} {\em \'{s}amatha} bd. \og Demeurer dans la paix,
développement de la paix\fg, une des méditations ({\em dhy\={a}na}) du
H\={\i}nay\={a}na. La simplicité est la clé de cette pratique qui
consiste à entrer en relation avec sa respiration, ou avec sa marche,
sans préméditation, sans but. Elle permet aux névroses et au bavardage
mental d'apparaître clairement. C'est d'elle que la méditation du Zen
S\={o}t\={o}, dite {\em zazen}, s'est inspiré, et où elle revêt une
importance capitale (\og Le Zen c'est {\em zazen}.~» dit D\={o}gen
Zenji). \\

%{\dn \8{f}^^FytA} {\em \'{s}\={u}nyat\={a}} [{\em \'{s}\={u}nya} +
%sf. {\em t\={a}}] f. mah. Vacuité. \\

{\dn \399w\388wA} {\em \'{s}raddh\={a}} f. mah. \og La confiance\fg,
première des Cinq Catégories ({\em bala}) de la voie de
l'unification. Confiance en soi inébranlable qui naît de l'abandon à
la fois de l'espoir et du désespoir. On essaiera d'éviter la
traduction \og foi\fg, en raison du possible malentendu avec le
théisme. \\

{\dn \399wAvk} {\em \'{s}r\={a}vaka} [{\em \'{s}ru} + sf. {\em aka}]
bd. \og Auditeur\fg, disciple qui a réellement, directement, entendu les
enseignements du Bouddha ({\em buddha}). \\

{\dn \399wAvkyAn} {\em \'{s}r\={a}vakay\={a}na} [{\em
\'{s}r\={a}vaka} + {\em y\={a}na}] n. hin. \og Véhicule de l'auditeur\fg,
premier des deux {\em y\={a}na} du  H\={\i}nay\={a}na. Il est centré
sur la compréhension des Quatre Nobles Vérités  ({\em \={a}ryasatya}). \\

{\dn s\2G} {\em sa\d{m}gha} m. bd. \og Communauté\fg, au sens
restreint désigne la congrégation des moines, au sens large tous les
bouddhistes. La {\em sa\d{m}gha} constitue l'un des Trois Joyaux ({\em
triratna}). \\

{\dn s\2\9{v}Ets(y} {\em sa\d{m}v\d{r}tisatya} \ [pf. {\em sam} +
{\em v\d{r}ti} + {\em satya}] n. mah. \og Vérité en accord avec le
manifeste\fg, vérité conventionnelle, une des deux catégories que
distingue la logique du Mah\={a}y\={a}na pour classer les
objets de connaissance (l'autre est la vérité ultime, ou {\em
param\={a}rthasatya}). \og Une vérité conventionnelle est un objet
per\c{c}u de fa\c{c}on dualiste par une connaissance valide directe
qui le per\c{c}oit en fait. L'absence d'existence réelle d'un vase est
un exemple de vérité ultime, le vase un exemple de vérité
conventionnelle. Ceci est valable pour tous les phénomènes, tous sans
exception possèdent ces deux vérités. [...] Les vérités
conventionnelles sont de deux sortes: vraies et fausses. (Les [Tenants
du Milieu] Conséquentialistes [{\em pr\={a}sa\d{n}g\={\i}kam\={a}\-dhyamika}]
n'opèrent pas de distinction entre réalité ou fausseté sur ce plan;
toutes les vérités conventionnelles étant pour eux dépourvues
d'existence réelle il ne peut y avoir de relatif vrai. Par contre du
point de vue d'une conscience mondaine ordinaire, et de ce point de
vue exclusivement, les vérités conventionnelles sont divisées en
vraies et fausses.) Un exemple de convention vraie est l'eau, un
exemple de convention fausse un mirage.~» (G. Khensur) \\

{\dn s\2sAr} {\em sa\.{m}s\={a}ra} [vr. {\em sa\.{m}s\d{r}}]
m. bd. Dans le livre III du Milindapa\~{n}ha ({\em Les questions de
Ménandre}) le roi gréco-bouddhique Ménandre interroge le moine
N\={a}gasena: \\
\og --- N\={a}gasena, tu parles du {\em sa\.{m}s\={a}ra}. Qu'est-ce que le
{\em sa\.{m}s\={a}ra}? \\
--- Un être naît sur cette terre et y meurt; mort ici, il renaît
ailleurs et y meurt, etc. Voilà ce qu'est le {\em sa\.{m}s\={a}ra}. \\
--- Donne-moi une comparaison. \\
--- Un homme mange une mangue et plante le noyau; ce noyau croît et un
grand manguier qui porte des fruits; un homme mange un de ces fruits
et plante le noyau, d'où croît un manguier, etc. Le point de départ de
cet enchaînement est inconnaissable. Il en est de même du
{\em sa\.{m}s\={a}ra}.\fg \\

{\dn s\2-kAr} {\em sa\.{m}sk\={a}ra} [vr. {\em sa\.{m}sk\d{r}}]
m. bd. \og Les formations mentales\fg, quatrième des Cinq Agrégats ({\em
skandha}). C'est l'intellect au sens de ce qui étiquette les choses,
leur donne un nom et les met dans certaines catégories
(concepts). $\mid$ bd. Seconde des douze causes de la
souffrance. cf. {\em prat\={\i}tyasamutp\={a}da}. \\

{\dn smy} {\em samaya} [{\em sami}] m. vaj. \og V{\oe}u\fg, engagement
du pratiquant qui se lie ainsi au {\em guru}, à la voie du bouddhisme
tantrique. \\

%{\dn smAED} {\em sam\={a}dhi} [{\em sam\={a}dhi}] m. bd. \\

{\dn sM\3E2wA} {\em samj\~{n}\={a}} f. bd. \og La perception\fg, troisième
des Cinq Agrégats ({\em skandha}). Elle est une forme aigu\"{e} de la
conscience sensorielle ({\em vedan\={a}}), orientée vers la
recherche des possibilités de modifier en apparence les projections
pour détourner les situations à notre profit. (C. Trungpa) \\

{\dn sADn} {\em s\={a}dhana} [{\em s\={a}dh} + sf. {\em ana}]
n. vaj. \og Visualisation\fg, pratique tantrique consistant à visualiser
une déité tutélaire désignée par le {\em guru} (en tibétain elle est
appelée {\em yidam}). Cette déité, ses attributs et sa forme
(irrité, semi-irrité ou apaisé) représente un aspect névrotique du
méditant qui devra en fusionnant avec elle réaliser sa propre vacuité
d'existence intrinsèque ainsi que celle de la déité, et ainsi faire
naître un aper\c{c}u de l'Éveil (cela revient à faire naître l'aspect
éveillé de la même déité). cf. {\em mah\={a}mudr\={a}}. \\

{\dn \8{s}/} {\em s\={u}tra} [{\em siv} + sf. {\em tra}] n. bd \og fil\fg,
parole du Bouddha ({\em buddha}) recueillie par écrit, et qui fixent sa
doctrine ({\em dharma}). Un {\em s\={u}tra} est toujours une parole annotée du
Bouddha, mais au sens large il peut désigner quelques recueils écrits
par de grands maîtres, comme le \og Soûtra de l'Estrade du Sixième
Patriarche [Ch'an] Houei-neng\fg. Il faut noter qu'en contexte bouddhique
un {\em s\={u}tra} est souvent très long, contrairement à
l'hindouisme où il constitue une suite d'aphorismes (cf. les Yogas\={u}tra de
Pata\~{n}jali). On préfèrera généralement la traduction \og soûtra~» qui
est acceptée par Le Petit Robert. Sinon, \og parole\fg, \og dit~» ou
\og discours~» conviennent mieux que \og sermon\fg, qui est connoté par le
christianisme. cf. {\em tantra}. $\|$ p\={a}li {\em sutta}. \\

{\dn \8{s}/EpVk} {\em s\={u}trapi\d{t}aka} [{\em s\={u}tra} + {\em
pi\d{t}aka}] m. n. bd. \og Corbeille des soûtras\fg, section du canon
bouddhique ({\em tripi\d{t}aka}) recueillant les discours {\em
s\={u}tra} de point de doctrine du Bouddha ({\em buddha}). \\

{\dn \8{s}/yAn} {\em s\={u}tray\={a}na} \ [{\em s\={u}tra} +
{\em y\={a}na}] n. bd. Les trois premiers {\em y\={a}na}. Les deux
premiers (le {\em \'{s}ravakay\={a}na},
puis le {\em pratyekabuddhay\={a}na}) constituent le
Hin\={a}y\={a}na. Le Mah\={a}y\={a}na est le troisième.
Cette dénomination est utile lorsque l'on veut distinguer le
Vajray\={a}na (dit alors {\em tantray\={a}na}). \\

{\dn -k^^FD} {\em skandha} m. bd. \og Agrégat\fg, les Cinq
Agrégats sur la base desquels l'ego est imputé et désigné. Dans
l'ordre du développement de l'ego, ce sont: la forme ({\em r\={u}pa}),
la sensation ({\em vedan\={a}}), la perception ({\em
sa\d{m}j\~{n}\={a}}), les formations mentales ({\em
sa\.{m}sk\={a}ra}) et la conscience ({\em
vij\~{n}\={a}na}). L'Ak\'{s}ayamatinirde\'{s}a dit: \og Les formes
s'apparentent à l'écume: comme l'écume elles ne sont
pas un soi, un être vivant, un principe vital, un principe nourricier,
un individu, une personne. Ce qu'est la nature de l'écume, cela est
aussi la nature des formes. [...] Les sensations s'apparentent à une
bulle, les formations à un [tronc de] bananier, les consciences à une
illusion. L'illusion n'est pas un soi, un être vivant, etc. Ce qu'est
la nature de l'illusion, cela est la nature de la conscience. [...]
Les agrégats s'apparentent à un rêve [ou à une émanation], et le rêve
n'est pas un soi, etc. Ce qu'est la nature du rêve, cela est aussi la
nature des agrégats. [...] Ce qu'on appelle agrégat est du monde, et
le monde a pour caractéristique la destruction. La nature du monde est
d'être naturellement impermanent ({\em anitya}), naturellement
souffrance ({\em du\d{h}kha}), vide  [cf. {\em \'{s}\={u}nyat\={a}}],
non-soi [{\em an\={a}tman}], paix. Qui est compétent sur ce point est
un héros pour l'Éveil [{\em bodhisattva}] dit habile vis-à-vis des
agrégats.~» cf. {\em an\={a}tman} $\|$ pâli {\em khanda} \\

{\dn -pf\0} {\em spar\'{s}a} [{\em sp\d{r}\'{s}} {\bf I}] m. bd. \og Le
toucher\fg, la prise de contact sensorielle (sixième des douze cause de
la souffrance. cf. {\em prat\={\i}tyasamutp\={a}da}). \\

{\dn -\9{m}Et} {\em sm\d{r}ti} [{\em sm\d{r}} + sf. {\em ti}]
f. mah. \og Conscience intelligente, attention\fg, troisième des Cinq
Catégories ({\em bala}) de la voie de l'unification. Le
pratiquant effectue un retour à sa vie quotidienne, à ses modes de
comportement mais il prête maintenant attention à ses objets
mentaux ainsi qu'à un certain sentiment d'être, sans réserve ni
distraction. \\

{\dn -\9{m}Etup-TAn} {\em sm\d{r}tiupasth\={a}na} [{\em sm\d{r}ti} +
{\em upa} + {\em sth\={a}na}] n. hin. \og Les quatre fondements de
l'attention\fg. Il s'agit de l'attention au corps, à la vie, à
l'énergie et à l'esprit, qui constituent le fondement de la méditation
{\em vipa\'{s}yan\={a}}. $\|$ p\={a}li {\em satipatthana}. \\

%{\dn hFnyAn} {\em h\={\i}nay\={a}na} [{\em h\={\i}na} + {\em
%y\={a}na}] n. mah. \og Petit Véhicule\fg \\

{\dn h\?zk} {\em heruka} vaj. Déité masculine irritée visualisée dans
les pratiques tantriques ({\em s\={a}dhana}). \\

{\dn h\?v\385w} {\em hevajra} np. vaj. Un {\em ma\d{n}\d{d}ala} de
l'{\em anuttarayogay\={a}na}. \\

%Outre les neufs {\em y\={a}na}, on peut aussi présenter le voyage spirituel
%bouddhique sous l'aspect des Cinq voies: accumulation, unification,
%vue, méditation et fin de l'apprentissage.

\end{document}
